% MAREA — validation only: design, results of the 10 m run and the planned
% pressure-sensor campaign. Standalone. Every table row is written by
%   python -m experiments.v10_validation_report --root results/server_10m \
%          --dutch-res 10 --granadeiro-status "under review"
% into docs/paper/tables/*_rows.tex (numbers.json holds every figure the text
% quotes), from the 10 m validation run on the group server (2026-10-01..06):
% notebooks (MAREA, no delay), experiments/v7 (NIDEM), v8 (Granadeiro,
% under review), v9 (scores, reproducing every notebook to 1e-16).
% Compile from docs/paper:  pdflatex validation_summary.tex (twice)
\documentclass[11pt,a4paper]{article}

\usepackage[english]{babel}
\usepackage[a4paper,top=2.2cm,bottom=2.2cm,left=2.5cm,right=2.5cm]{geometry}
\usepackage{amsmath}
\usepackage{booktabs}
\usepackage{graphicx}
\usepackage{microtype}
\usepackage{float}
\usepackage[colorlinks=true,allcolors=blue]{hyperref}

\setlength{\parskip}{0.4em}
\setlength{\parindent}{0pt}
% table rows written by experiments/v10_validation_report.py; \input is not
% safe inside a tabular, the TeX primitive is
\makeatletter\newcommand\inputrows[1]{\@@input #1 }\makeatother

\title{MAREA: validation design, results to date\\and the pressure-sensor campaign}
\author{Jorge Al\'ias (AIC, Universidad de Oviedo) \and Fernando Garc\'ia-Gonz\'alez (IEO-CSIC) \and Pablo Gonz\'alez (AIC, Universidad de Oviedo)}
\date{Working document, 6 October 2026}

\begin{document}
\maketitle

\section{What is validated, and why}

MAREA produces two quantities, and each one is checked separately.
\begin{itemize}
\item \textbf{The interior tidal lag} $\tau$ of each 3-km section of the
  estuary: the delay with which the tide reaches that section after
  reaching the mouth. Validating it answers two questions: is the lag
  read from the archive the real one, and does applying it improve the
  water level assigned inside the estuary?
\item \textbf{The elevation} $z$ of each intertidal pixel, from which the
  emerged time follows. Validating it gives the error of the map and how
  that error grows inland, where the lag matters most.
\end{itemize}
Each quantity is tested on site, in the R\'ia de Villaviciosa, and
externally, against records and surveys made by others
(Table~\ref{tab:blocks}). Nothing that judges a result enters a fit: the
inner tide gauges are never used as boundary, and a quarter of the
Villaviciosa RTK-GNSS points (91 of 361) was sealed as a test set before
any elevation product existed.

\begin{table}[H]
\centering
\caption{The validation blocks.}
\label{tab:blocks}
\small
\setlength{\tabcolsep}{4pt}
\resizebox{\textwidth}{!}{%
\begin{tabular}{@{}lllll@{}}
\toprule
block & reference & sites & metric & status \\
\midrule
lag, on site & pressure sensors along the r\'ia & Villaviciosa & error of the level difference & planned \\
 & & & mouth--station, $E(\tau)$ & (Section~\ref{sec:pressure}) \\
lag, external & tide-gauge pairs: mouth and & Ferrol, Westerschelde, & same as above & 10~m run \\
 & inner gauge held out & Ems-Dollard, Wadden & & \\
elevation, on site & RTK-GNSS survey (points) & Villaviciosa & RMSE, slope $b$, $r$, bias & 10~m run \\
elevation, on site & IGN PNOA LiDAR, 2~m (whole map) & Villaviciosa & RMSE, slope $b$, $r$, bias & 10~m run \\
elevation, external & Vaklodingen bathymetric & Westerschelde, Ems-Dollard, & RMSE, slope $b$, $r$ & 10~m run \\
 & and intertidal survey (NL) & Wadden Sea & & \\
\bottomrule
\end{tabular}}
\end{table}

\section{Metrics, datums and methods compared}

\paragraph{Lag.} Between a gauge at the mouth (level $h_M$) and a gauge
inside ($h_I$), each centred on its own mean over 2023--2025, the
measured level difference is $m(t) = h_M(t) - h_I(t)$. A lag $\tau$
predicts it as $p_\tau(t) = h_M(t) - h_M(t-\tau)$: the inner level is the
mouth's, $\tau$ minutes later. The error of that prediction,
\begin{equation}
E(\tau) = \sqrt{\frac{1}{N}\sum_{t}\big(m(t) - p_\tau(t)\big)^2}
        = \sqrt{\frac{1}{N}\sum_{t}\big(h_M(t-\tau) - h_I(t)\big)^2},
\label{eq:lag}
\end{equation}
over every 10-min instant of 2023--2025 with both records, is reported for
no lag ($\tau = 0$: one uniform level over the estuary), for the lag MAREA
applies at the inner gauge, and for the best constant lag, which no lag
can beat. The slope $b$ and correlation $r$ of $p_\tau$ on $m$ say
whether MAREA's lag gives the right size ($b = 1$; below 1 the lag is too
short, above 1 too long) and shape of the difference. The lag between the
gauges in reaching the same level (mean over paired tidal cycles, both
limbs, three levels) is given for reference.

\paragraph{Elevation.} Product and truth are on different vertical datums,
so each series is centred on its own median, $y_i = z_i -
\operatorname{median}(z)$, $x_i = z^{\mathrm{t}}_i -
\operatorname{median}(z^{\mathrm{t}})$, and the removed offset is reported
as the bias. With the truth on the abscissa,
\begin{equation}
\mathrm{RMSE} = \sqrt{\frac{1}{n}\sum_{i}(y_i - x_i)^2},\qquad
b = \frac{\sum_i (x_i-\bar x)(y_i-\bar y)}{\sum_i (x_i-\bar x)^2},\qquad
r = \frac{\sum_i (x_i-\bar x)(y_i-\bar y)}{\sqrt{\sum_i (x_i-\bar x)^2\sum_i (y_i-\bar y)^2}}.
\label{eq:elev}
\end{equation}
The slope $b$ says whether the map reproduces the relief ($b = 1$) or
compresses it ($b < 1$); $r$, how much of it the map explains. Only truth
within the range of tide levels the archive sampled ($\pm 0.25$~m) is
scored.

\paragraph{Datums.} The maps are referenced to the boundary tide: the EOT20
ocean tide model, whose zero is the mean sea surface. The Villaviciosa
RTK-GNSS survey and the PNOA LiDAR are both in the Spanish national datum
(REDNAP, mean sea level at Alicante): the survey's ellipsoidal heights
were converted with the IGN geoid EGM08-REDNAP ($N \approx 52.74$~m in
Villaviciosa). In that common datum the LiDAR and the RTK agree to a
median difference of 0.00~m (centred RMSE 0.12~m, $r = 0.98$, on the 270
development points), so the two are interchangeable references. Because
every metric is median-centred, the datum changes only the bias, which in
Villaviciosa measures where the maps' zero lies against the national
datum. Tide gauges are compared after removing their own mean; the
Vaklodingen (NAP) enter only through centred metrics.

\paragraph{Methods compared.} Every method works on the same archive
(Sentinel-2 L2A cube of the site, 2023--2025, 10~m grid at every site, the
same boundary tide) and is scored on the same pixels (the intertidal
record of the MAREA extraction), against the same truth, with the same
metrics. On top of that, each method applies its own published rules for
selecting scenes and masking pixels:
\begin{itemize}
\item \textbf{No delay}: MAREA's elevation inversion with the lag set to
  zero, on MAREA's record. It isolates what the lag contributes.
\item \textbf{MAREA}: the same inversion with the interior lag. Its record
  keeps a scene when the whole frame is at most 5\,\% cloudy or, failing
  that, when the transition zone (the pixels that change between water and
  land) is at most 10\,\% cloudy; clouds inside kept scenes are masked per
  pixel.
\item \textbf{NIDEM} (Bishop-Taylor et al., 2019; ITEM, Sagar et al.,
  2017): waterline-contour interpolation over tide intervals with ITEM's
  all-observations rule and pixel-quality mask.
\item \textbf{Granadeiro et al.\ (2021)}, faithful reimplementation (their
  logistic fit with \texttt{nplr}, lag sample, \texttt{mgcv} GAM and scene
  rule, $<10\,\%$ cloud over the frame; Sen2Cor instead of ACOLITE). It has
  been run at every site, but its outputs show signs of an input problem
  (an NDWI dispersion above its theoretical maximum of 1 in the intertidal
  mask, near-constant preliminary elevations); its scores are withheld
  until that is resolved.
\end{itemize}
Each method is scored on every pixel it resolves (\emph{all}) and on the
pixels every reported method resolves (\emph{common}), the fair
head-to-head.

\section{Results of the 10 m run}

\subsection{Elevation on site: RTK-GNSS survey, R\'ia de Villaviciosa}

The survey holds 361 fixed points taken at the lowest spring low water of
the month; the 270 development points are scored here (194 fall on the
mapped intertidal pixels), the test points stay sealed.

\begin{table}[H]
\centering
\caption{Elevation against the RTK-GNSS survey, Villaviciosa development
subset, heights in the Alicante datum. Scenes: number each method's own
rule keeps; $n$: surveyed points the map resolves (the same 194 for every
method); RMSE, $b$, $r$ as in Eq.~\eqref{eq:elev}; bias: median of map
minus survey, the position of the maps' zero against the national datum.}
\label{tab:rtk}
\small
\begin{tabular}{@{}lrrrrrr@{}}
\toprule
method & scenes & $n$ & RMSE (m) & $b$ & $r$ & bias (m) \\
\midrule
\inputrows{tables/rtk_rows.tex}
\bottomrule
\end{tabular}
\end{table}

MAREA and its no-delay twin are the most accurate maps on the surveyed
flats (0.19~m), NIDEM follows at 0.23~m, although it reproduces the relief
better ($b = 0.85$ against 0.73). The maps' zero lies 0.08~m above the
Alicante datum. The lag applied on the surveyed pixels is a few minutes:
Villaviciosa tests the elevation estimator more than the lag, which is why
the on-site lag campaign of Section~\ref{sec:pressure} is needed.

\subsection{Elevation on site: LiDAR over the whole map, R\'ia de Villaviciosa}

The IGN PNOA LiDAR DTM (2~m, second coverage, Alicante datum) covers the
whole r\'ia. A LiDAR returns the water surface where water stood during
the flight, so water is removed at 2~m first: a 6-m window flatter than
5~mm is a water surface (7\,\% of the pixels, mostly two flight strips with
water at $-0.97$ and $-0.57$~m). Each 10~m cell is the mean of its ground
pixels, kept when ground covers at least 80\,\% of the cell. A one-pixel
shift of the LiDAR raises MAREA's error in every direction, so the two
grids are co-registered.

\begin{table}[H]
\centering
\caption{Elevation against the PNOA LiDAR DTM, Villaviciosa, on 10~m cells.
$n$, RMSE, $b$, $r$ on every cell each map resolves; RMSE common on the
19\,198 cells every method resolves; bias against the Alicante datum.}
\label{tab:lidar}
\small
\begin{tabular}{@{}lrrrrrr@{}}
\toprule
method & $n$ & RMSE (m) & $b$ & $r$ & RMSE common (m) & bias (m) \\
\midrule
\inputrows{tables/lidar_rows.tex}
\bottomrule
\end{tabular}
\end{table}

On the cells every method resolves, the three maps agree with the LiDAR to
within 1~cm of each other (0.29--0.30~m), with NIDEM again closer to the
full relief ($b = 0.81$). MAREA and no delay map about 6\,000 cells more
than NIDEM, the lowest ground beside the channels; those cells carry most
of their higher error over all cells (0.45~m), and are where a per-pixel
quality filter for MAREA is needed. Near the RTK points the maps and the
LiDAR agree to about 0.12~m: the surveyed flats are the easier part of the
map.

\subsection{Lag against tide-gauge pairs}

\begin{table}[H]
\centering
\caption{Interior lag against tide-gauge pairs, 2023--2025. Scenes:
MAREA's record. Lag between the gauges in reaching the same level
(positive: inner gauge later); lag MAREA applies at the inner gauge, with
in brackets the 95\,\% bootstrap interval of its 3-km section (30
resamplings of the scenes). $E$: error of the predicted level difference
between mouth and inner gauge, Eq.~\eqref{eq:lag}, without lag, with
MAREA's lag and with the best constant lag (in brackets, that lag). $b$,
$r$: slope and correlation of MAREA's predicted difference on the measured
one ($b < 1$: lag too short; $b > 1$: too long).}
\label{tab:lag}
\small
\setlength{\tabcolsep}{3.5pt}
\resizebox{\textwidth}{!}{%
\begin{tabular}{@{}lrrrrrrrr@{}}
\toprule
 & & lag, gauges & lag, MAREA (min) & $E$ (m) & $E$ (m) & $E$ (m), best & & \\
site (inner gauge, km) & scenes & (min) & [95\,\% CI of band] & no lag & MAREA & (lag, min) & $b$ & $r$ \\
\midrule
\inputrows{tables/lag_rows.tex}
\bottomrule
\end{tabular}}
\end{table}

Where the tide reaches the inner gauge late (Westerschelde, Ems-Dollard,
Wadden), MAREA's lag removes most of the error of a uniform level: from
0.28 to 0.20~m, 0.50 to 0.26~m and 0.53 to 0.23~m, that is 48, 73 and
81\,\% of it. In the Ems-Dollard it is close to the best any constant lag
achieves (0.26 against 0.25~m). In the Westerschelde the lag is too short
(8 against 20~min, $b = 0.30$) and in the Wadden too long (121 against
93~min, $b = 1.21$). In Ferrol, with no lag to measure (4~min), MAREA
applies none.

\subsection{Elevation against the Vaklodingen}

\begin{table}[H]
\centering
\caption{Elevation against the Rijkswaterstaat Vaklodingen, the official
bathymetric and intertidal survey of the Dutch coast (20~m cells, NAP),
Eq.~\eqref{eq:elev}. Maps on a 10~m grid. Scenes: number each method's own
rule keeps. $n$, RMSE, $b$, $r$ on every intertidal pixel each map
resolves; last column, RMSE on the pixels every reported method resolves
(common subset, $n_c$). $^{b}$~Not computable with its own rule: the
lowest tide interval holds one scene, valid on almost no pixel.}
\label{tab:external}
\small
\setlength{\tabcolsep}{4pt}
\begin{tabular}{@{}lrrrrrr@{}}
\toprule
method & scenes & $n$ & RMSE (m) & $b$ & $r$ & RMSE common (m) \\
\midrule
\inputrows{tables/external_rows.tex}
\bottomrule
\end{tabular}
\end{table}

In the Ems-Dollard, where the tide reaches the inner flats about an hour
late, applying the lag is what makes the map usable: the error falls from
0.60 to 0.35~m and the slope rises from 0.10 to 0.64. In the
Westerschelde, where the lag is small, MAREA and no delay are equivalent
(0.40--0.41~m) and NIDEM is slightly more accurate (0.37~m). In the
Wadden, MAREA's map is worse than no delay (0.41 against 0.37~m on all
pixels, 0.38 against 0.37~m on the common ones) and NIDEM is the most
accurate (0.29~m): the lag MAREA applies there overshoots the gauges' by
about half an hour (Table~\ref{tab:lag}), and the reason this degrades the
elevation while it improves the water level is under investigation.

\section{Further work: the pressure-sensor campaign in Villaviciosa}\label{sec:pressure}

The R\'ia de Villaviciosa has no tide gauge inside, so the lag MAREA reads
there has never been checked against a measurement. A short campaign of
pressure sensors along the r\'ia measures it directly and closes the
on-site lag block. The MAREA lags to be tested already exist, so the
campaign is a prediction checked afterwards, not a calibration.

\paragraph{Stations.} One sensor per section of the product, plus the
mouth as reference (Table~\ref{tab:stations}). Each is fixed to a stake
or mooring at the channel margin, as low as the channel allows, so that
it stays submerged through as much of the tidal cycle as possible; where
the channel dries at low water, lags are computed only for levels above
the sensor.

\begin{table}[H]
\centering
\caption{Planned stations and the MAREA lag each one will test
(Villaviciosa product, 2023--2025 record of 97 scenes; sections of 3~km
from the mouth).}
\label{tab:stations}
\small
\begin{tabular}{@{}llrr@{}}
\toprule
station & position & from mouth (km) & MAREA lag (min) [95\,\% interval] \\
\midrule
\inputrows{tables/stations_rows.tex}
P3 (optional) & inner end of the flats & $\approx 9$ & extrapolated from section 2 \\
\bottomrule
\end{tabular}
\end{table}

\paragraph{Instruments and deployment.}
\begin{itemize}
\item Absolute pressure loggers at every station and a barometer on land
  (or a logger kept in air) for atmospheric compensation, all sampling at
  1~min (5~min at most) and synchronised to UTC before deployment; clock
  drift is read at recovery and corrected linearly.
\item At least 29 days, one full spring--neap cycle, so that every level
  and both limbs are sampled over the range the archive saw.
\item Each sensor's elevation $z_s$ surveyed by RTK-GNSS at deployment and
  at recovery, in the same frame as the elevation survey.
\item A conductivity reading (CTD cast or conductivity logger) at each
  station at deployment and recovery, because brackish water changes the
  density $\rho$ and therefore the level scale by up to 2--3\,\%. This
  affects the amplitude of the record, not its timing.
\end{itemize}

\paragraph{Processing.} Pressure is converted to water level with
\begin{equation}
h = z_s + \frac{p - p_{\mathrm{atm}}}{\rho\, g},
\end{equation}
with $p$ the logger pressure, $p_{\mathrm{atm}}$ the barometer record and
$\rho$ from the measured salinity (1025~kg\,m$^{-3}$ for sea water). Each
record is put on a common 1-min grid, spikes removed, gaps longer than
one hour left open, and demeaned. The code already exists in the library
(module \texttt{lag\_validation} of \texttt{pyintertidal}); the same
functions produced Table~\ref{tab:lag}.

\paragraph{Analysis.} Every station is set against P0 exactly as the
gauge pairs of Table~\ref{tab:lag}:
\begin{enumerate}
\item the error $E(\tau)$ of the level difference with the station,
  Eq.~\eqref{eq:lag}, without lag, with MAREA's lag and with the best
  constant lag, and the slope $b$ and correlation $r$;
\item the lag in reaching the same level, per limb, at the 25th, 50th and
  75th percentile of P0's level, compared with the MAREA lag of
  Table~\ref{tab:stations};
\item at the Sentinel-2 overpasses that fall inside the campaign, the
  measured local water level against the one MAREA assigns (boundary plus
  lag): the exact quantity the elevation inversion uses.
\end{enumerate}

\paragraph{What counts as success, fixed in advance.} The measured lag at
each station falls inside the 95\,\% interval of its section, and MAREA's
lag lowers $E$ at the inner stations. A measured lag outside the interval
is reported as such and fed back as a known limitation, not tuned away.

\paragraph{External elevation references.} Beyond Villaviciosa, the
elevation will be tested against bathymetric LiDAR or official
bathymetries wherever they cover intertidal ground at low water, to
establish the error and its growth inland across regions and tidal
regimes. The Dutch Vaklodingen already serves this purpose
(Table~\ref{tab:external}).

\end{document}
